\documentclass[10pt,a4paper]{amsart}
\usepackage{amsmath,amssymb,amsthm,graphicx}
\usepackage{verbatim}
\newenvironment{reference}{\comment}{\endcomment}
%\newenvironment{reference}{}{}
\newenvironment{slogan}{\comment}{\endcomment}
\newenvironment{history}{\comment}{\endcomment}
% For commutative diagrams we use Xy-pic
\usepackage[all]{xy}

% We use 2cell for 2-commutative diagrams.
\xyoption{2cell}
% Theorem environments.
%
\theoremstyle{plain}
\newtheorem{theorem}[subsection]{Theorem}
\newtheorem{proposition}[subsection]{Proposition}
\newtheorem{lemma}[subsection]{Lemma}

\theoremstyle{definition}
\newtheorem{definition}[subsection]{Definition}
\newtheorem{example}[subsection]{Example}
\newtheorem{exercise}[subsection]{Exercise}
\newtheorem{situation}[subsection]{Situation}

\theoremstyle{remark}
\newtheorem{remark}[subsection]{Remark}
\newtheorem{remarks}[subsection]{Remarks}

\numberwithin{equation}{subsection}

% Macros
%
\def\lim{\mathop{\mathrm{lim}}\nolimits}
\def\colim{\mathop{\mathrm{colim}}\nolimits}
\def\Spec{\mathop{\mathrm{Spec}}}
\def\Hom{\mathop{\mathrm{Hom}}\nolimits}
\def\Ext{\mathop{\mathrm{Ext}}\nolimits}
\def\SheafHom{\mathop{\mathcal{H}\!\mathit{om}}\nolimits}
\def\SheafExt{\mathop{\mathcal{E}\!\mathit{xt}}\nolimits}
\def\Sch{\mathit{Sch}}
\def\Mor{\mathop{\mathrm{Mor}}\nolimits}
\def\Ob{\mathop{\mathrm{Ob}}\nolimits}
\def\Sh{\mathop{\mathit{Sh}}\nolimits}
\def\NL{\mathop{N\!L}\nolimits}
\def\CH{\mathop{\mathrm{CH}}\nolimits}
\def\proetale{{pro\text{-}\acute{e}tale}}
\def\etale{{\acute{e}tale}}
\def\QCoh{\mathit{QCoh}}
\def\Ker{\mathop{\mathrm{Ker}}}
\def\Im{\mathop{\mathrm{Im}}}
\def\Coker{\mathop{\mathrm{Coker}}}
\def\Coim{\mathop{\mathrm{Coim}}}

% Boxtimes
%
\DeclareMathSymbol{\boxtimes}{\mathbin}{AMSa}{"02}

%
% Macros for moduli stacks/spaces
%
\def\QCohstack{\mathcal{QC}\!\mathit{oh}}
\def\Cohstack{\mathcal{C}\!\mathit{oh}}
\def\Spacesstack{\mathcal{S}\!\mathit{paces}}
\def\Quotfunctor{\mathrm{Quot}}
\def\Hilbfunctor{\mathrm{Hilb}}
\def\Curvesstack{\mathcal{C}\!\mathit{urves}}
\def\Polarizedstack{\mathcal{P}\!\mathit{olarized}}
\def\Complexesstack{\mathcal{C}\!\mathit{omplexes}}
% \Pic is the operator that assigns to X its picard group, usage \Pic(X)
% \Picardstack_{X/B} denotes the Picard stack of X over B
% \Picardfunctor_{X/B} denotes the Picard functor of X over B
\def\Pic{\mathop{\mathrm{Pic}}\nolimits}
\def\Picardstack{\mathcal{P}\!\mathit{ic}}
\def\Picardfunctor{\mathrm{Pic}}
\def\Deformationcategory{\mathcal{D}\!\mathit{ef}}
\usepackage{cleveref}
\crefname{originallemma}{original lemma}{original lemmas}
\Crefname{originallemma}{Original lemma}{Original lemmas}
\makeatletter
\newlabel{lemma-stalk-push-finite}{{0DK3}{source}{}{sourcepointer-1}{}}
\newlabel{lemma-stalk-push-finite@cref}{{[originallemma][0][]0DK3}{[1][0][]source}{}{}{}}
\makeatother
\crefname{originallemma}{original lemma}{original lemmas}
\Crefname{originallemma}{Original lemma}{Original lemmas}
\makeatletter
\newlabel{spaces-properties-lemma-stalk-pullback}{{04K2}{source}{}{sourcepointer-2}{}}
\newlabel{spaces-properties-lemma-stalk-pullback@cref}{{[originallemma][0][]04K2}{[1][0][]source}{}{}{}}
\makeatother
\crefname{originaltheorem}{original theorem}{original theorems}
\Crefname{originaltheorem}{Original theorem}{Original theorems}
\makeatletter
\newlabel{spaces-properties-theorem-exactness-stalks}{{04K5}{source}{}{sourcepointer-3}{}}
\newlabel{spaces-properties-theorem-exactness-stalks@cref}{{[originaltheorem][0][]04K5}{[1][0][]source}{}{}{}}
\makeatother
\title{The projection formula for a finite morphism of algebraic spaces with sheaves of rings, checked on stalks}
\author{The Stacks Project Authors}
\date{stacks-project:spaces-cohomology.tex:local-snapshot-20261012; GNU FDL 1.2 or later}
\begin{document}
\maketitle
\noindent\textit{Editorial note.} All delivered mathematical text is the author's own, copied from the pinned Stacks Project chapter file spaces-cohomology.tex; no statement or proof step was rewritten, reordered or removed.\par\medskip
\noindent\textit{Editorial note.} Preamble provenance: the pinned chapter file opens with \textbackslash{}input\{preamble\} and carries no preamble text of its own, so the project's single Stacks-wide preamble staged at sources/stacks-spaces-cohomology/preamble.tex is excerpted verbatim as three spans, lines 8-12 (verbatim package and the project's own reference/slogan/history display environments), lines 14-18 (Xy-pic setup) and lines 154-217 (theorem declarations and mathematical macros). The wrapper class and the amsmath/amssymb/amsthm/graphicx base packages are supplied by the assembler. No macro was invented or redefined.\par\medskip
\noindent\textit{Editorial note.} Same-chapter pointers are located at their \textbackslash{}label lines inside the pinned chapter file; cross-chapter pointers are located in the project's own tag list staged in the snapshot (sources/stacks-chow/tags.txt), at the verbatim lines carrying their canonical tags.\par\medskip
\noindent\textit{Allowed standard prerequisites (editorial).}\begin{itemize}
\item A morphism of algebraic spaces is an isomorphism exactly when it is so on stalks at geometric points. Properties of Spaces, theorem spaces-properties-theorem-exactness-stalks (invoked at proof line 578). It reduces the verification of the isomorphism of module sheaves to the corresponding statement about stalks at geometric points.
\item The behaviour of stalks under pullback and under pushforward along a finite morphism. Properties of Spaces, label spaces-properties-lemma-stalk-pullback (invoked at proof line 600) and Same chapter spaces-cohomology.tex, label lemma-stalk-push-finite (invoked at proof line 602). The first describes the stalks of \$f\textasciicircum{}*\textbackslash{}mathcal\{G\}\$, the second identifies the stalk of a pushforward along a finite morphism with the product of the stalks over the finitely many geometric lifts.
\item Tensor products commute with taking stalks. Standard, used at proof lines 597-599. It allows the two sides of the displayed identity of stalks to be rewritten as tensor products over the residue field rings, where the equality is formal.
\end{itemize}\medskip
\ifdefined\hypertarget\hypertarget{sourcepointer-1}{}\fi
\noindent\textit{Source reference.} Original lemma 0DK3: sources/stacks-spaces-cohomology/pinned-original-source.tex, source lines 518--518. Navigation pointer for \textbackslash{}ref\{lemma-stalk-push-finite\} used inside the counted span (line 602); the argument that consumes it is retained in the span.\par\medskip
\ifdefined\hypertarget\hypertarget{sourcepointer-2}{}\fi
\noindent\textit{Source reference.} Original lemma 04K2: sources/stacks-chow/tags.txt, source lines 5958--5958. Navigation pointer for \textbackslash{}ref\{spaces-properties-lemma-stalk-pullback\} used inside the counted span (line 600); the argument that consumes it is retained in the span.\par\medskip
\ifdefined\hypertarget\hypertarget{sourcepointer-3}{}\fi
\noindent\textit{Source reference.} Original theorem 04K5: sources/stacks-chow/tags.txt, source lines 5961--5961. Navigation pointer for \textbackslash{}ref\{spaces-properties-theorem-exactness-stalks\} used inside the counted span (line 579); the argument that consumes it is retained in the span.\par\medskip
\begin{lemma}
\label{lemma-finite-rings}
Let $S$ be a scheme. Let $\pi : X \to Y$ be a finite morphism of algebraic
spaces over $S$. Let $\mathcal{A}$ be a sheaf of rings on $X_\etale$.
Let $\mathcal{B}$ be a sheaf of rings on $Y_\etale$.
Let $\varphi : \mathcal{B} \to \pi_*\mathcal{A}$
be a homomorphism of sheaves of rings so that we obtain a
morphism of ringed topoi
$$
f = (\pi, \varphi) :
(\Sh(X_\etale), \mathcal{A})
\longrightarrow
(\Sh(Y_\etale), \mathcal{B}).
$$
For a sheaf of $\mathcal{A}$-modules $\mathcal{F}$ and a
sheaf of $\mathcal{B}$-modules $\mathcal{G}$ the canonical map
$$
\mathcal{G} \otimes_\mathcal{B} f_*\mathcal{F}
\longrightarrow
f_*(f^*\mathcal{G} \otimes_\mathcal{A} \mathcal{F}).
$$
is an isomorphism.
\end{lemma}

\begin{proof}
The map is the map adjoint to the map
$$
f^*\mathcal{G} \otimes_\mathcal{A}
f^* f_*\mathcal{F} =
f^*(\mathcal{G} \otimes_\mathcal{B} f_*\mathcal{F})
\longrightarrow
f^*\mathcal{G} \otimes_\mathcal{A} \mathcal{F}
$$
coming from $\text{id} : f^*\mathcal{G} \to f^*\mathcal{G}$
and the adjunction map $f^* f_*\mathcal{F} \to \mathcal{F}$.
To see this map is an isomorphism, we may check on stalks
(Properties of Spaces, Theorem
\ref{spaces-properties-theorem-exactness-stalks}).
Let $\overline{y}$ be a geometric point of $Y$ and
let $\overline{x}_1, \ldots, \overline{x}_n$ be the geometric
points of $X$ lying over $\overline{y}$.
Working out what our maps does on stalks, we see that we
have to show
$$
\mathcal{G}_{\overline{y}}
\otimes_{\mathcal{B}_{\overline{y}}}
\left(
\bigoplus\nolimits_{i = 1, \ldots, n} \mathcal{F}_{\overline{x}_i}
\right) =
\bigoplus\nolimits_{i = 1, \ldots, n}
(\mathcal{G}_{\overline{y}}
\otimes_{\mathcal{B}_{\overline{x}}}
\mathcal{A}_{\overline{x}_i}) \otimes_{\mathcal{A}_{\overline{x}_i}}
\mathcal{F}_{\overline{x}_i}
$$
which holds true. Here we have used that
taking tensor products commutes with taking stalks, the
behaviour of stalks under pullback
Properties of Spaces, Lemma \ref{spaces-properties-lemma-stalk-pullback}, and
the behaviour of stalks under pushforward along a closed immersion
Lemma \ref{lemma-stalk-push-finite}.
\end{proof}

\end{document}
